\documentclass{article}
\usepackage{amsmath,amssymb,amsthm}
\newtheorem{theorem}{Theorem}
\title{A compact Cauchy obstruction to a common amplitude pool}
\date{24 November 2026}
\begin{document}
\maketitle

\paragraph{Status and scope.}
The obstruction below is kernel checked in
\texttt{HeadComplexity/Results/LargeDomainCauchySharedPool.lean}.
It concerns one proposed homogeneous, zero-bias induction for a growing
permutation decoder. Independent node biases are not included in its
inequalities and can evade this particular three-node obstruction.
It does not rule out target-adaptive residues, a full joint decoder,
or any upper bound for the original quantity $W(n)$.
No original model definition is changed.

For $q\geq10$, let
\[
 K_{ij}=\frac1{(q-i-j)^2},\qquad 0\leq i,j\leq2,
 \qquad \sigma=(-1,1,1),\quad\tau=(1,-1,1).
\]
For a sign vector $v$, write $D_v$ for the corresponding diagonal matrix.

\begin{theorem}
There is no real vector $z$ such that both
$D_\sigma K D_\sigma z$ and $D_\tau K D_\tau z$ have all coordinates
strictly positive. This remains true for arbitrarily large $q$, hence
for arbitrarily compact normalized poles $0,1/q,2/q$.
\end{theorem}
\begin{proof}
Put
\[
 a=K_{00},\quad b=K_{01},\quad c=K_{02}=K_{11},
 \quad e=K_{12},\quad f=K_{22},\quad\Delta=cf-e^2.
\]
The two row-zero inequalities imply $az_0-bz_1>0$, and the two
row-one inequalities imply $-bz_0+cz_1>0$. Since
\[
 ac-b^2=\frac{2q^2-4q+1}{q^2(q-1)^4(q-2)^2}>0,
\]
multiplying the first by $c$ and the second by $b$ gives $z_0>0$.
Thus positivity of the first amplitude is forced by the two systems.
In fact their row-two sum also forces $z_2>0$, and the leading
principal inverse forces $z_1>0$.

All five entries are positive, and
\[
 \Delta=\frac{2q^2-12q+17}{(q-4)^2(q-3)^4(q-2)^2}>0.
\]
The proposed inequalities include the following three:
\[
 az_0-bz_1-cz_2>0,\qquad
 -bz_0+cz_1-ez_2>0,\qquad
 cz_0-ez_1+fz_2>0.
\]
Multiply them respectively by the positive quantities
$\Delta$, $bf+ce$, and $c^2+be$. Their sum is exactly
\[
 \bigl(a\Delta-b^2f+c^3\bigr)z_0.
\]
But
\[
 a\Delta-b^2f+c^3=
 \frac{-2P(q)}{q^2(q-4)^2(q-3)^4(q-2)^6(q-1)^4}<0,
\]
where
\[
 P(q)=2q^{10}-40q^9+339q^8-1584q^7+4432q^6-7488q^5
       +7135q^4-2808q^3-756q^2+912q-136.
\]
To verify the last sign without an asymptotic argument, put $r=q-10\geq0$.
Then
\[
\begin{split}
 P(q)={}&2r^{10}+160r^9+5739r^8+121536r^7+1682752r^6
 +15916032r^5\\
 &+104140735r^4+465434592r^3+1359716004r^2
 +2344483392r+1811675384>0.
\end{split}
\]
Thus the same sum is strictly negative because $z_0>0$, a contradiction.
The original Lean declaration
\texttt{HeadComplexity.LargeDomain.no\_shared\_positive\_cauchy\_pool}
assumes $z_0>0$ and the three selected inequalities. The strengthened
\texttt{HeadComplexity.LargeDomain.no\_shared\_cauchy\_pool}
derives $z_0>0$ from five of the six rows and has no amplitude sign
assumption.
\end{proof}

\paragraph{Why separate one-level success is insufficient.}
The matrix $K$ is positive definite: it is the Gram matrix of the three
functions $e^{it}$ in the positive integral
\[
 K_{ij}=\int_0^\infty t e^{-qt}e^{it}e^{jt}\,dt.
\]
Linear independence of the exponentials gives strict definiteness.
Every positive definite real symmetric matrix $A$ admits some $z>0$
with $Az>0$. For example, minimize
$\frac12z^{\mathsf T}Az-\sum_i\log z_i$ on the positive orthant;
coercivity and the logarithmic barrier give an interior minimizer, whose
stationarity equation is $(Az)_i=1/z_i>0$.
Apply this separately to $D_\sigma K D_\sigma$ and
$D_\tau K D_\tau$. Thus each of the two sign problems has positive
amplitudes, but the amplitudes cannot be shared. This last explanatory
argument is a written proof, not part of the Lean declaration.

\paragraph{Consequence for the attempted induction.}
A one-level positive definite Cauchy argument can choose amplitudes after
seeing that level's signs. A recursive decoder cannot silently reuse those
amplitudes for a second independently specified sign problem. Already
three compact nodes give a counterexample to that intermediate lemma.
The theorem leaves open target-adaptive signed residues, independent
per-column amplitudes, different source gates, or cancellation involving
several levels at once.
\end{document}
